\chapter{Covariant Derivatives, Connections, and the Clean UBT Geometry}
\label{ch:covariant-tetrad}

This chapter explains the geometric core of UBT in the form closest to the
original intuition of the theory:
\[
 \Theta\quad\longrightarrow\quad E_\mu:=D_\mu\Theta
 \quad\longrightarrow\quad
 \{E_\mu,E_\nu\}\ \hbox{and}\ [D_\mu,D_\nu].
\]
The anticommutator carries the metric, while the commutator of covariant
derivatives carries curvature.  No embedding map, compact-fiber average, or
observer projection is needed in the minimal local construction.

\section{Ordinary, partial, and covariant derivatives}
For an ordinary scalar function $f(x)$, the derivative $df/dx$ measures how the
number $f$ changes when $x$ changes.  For a field on spacetime we use partial
derivatives,
\[
 \partial_\mu\Theta:=\frac{\partial\Theta}{\partial x^\mu},
 \qquad \mu=0,1,2,3.
\]
A partial derivative compares components written in a fixed basis.  This is
sufficient if the same basis can be used everywhere.  Geometry becomes subtler
when the local frame itself rotates or boosts from point to point.

A simple analogy is a vector written in a rotating coordinate frame.  Even a
physically constant vector acquires changing components because the basis is
changing.  The covariant derivative corrects for this basis change:
\[
 D_\mu\Theta=\partial_\mu\Theta+\rho_*(\Omega_\mu)\Theta.
\]
Here $\Omega_\mu$ is the connection and $\rho_*$ tells us how the connection
acts in the representation of $\Theta$.  For biquaternions, a left action, a
right action, and a two-sided action are genuinely different and must not be
silently interchanged.

\section{What is a connection?}
A connection is not a force by itself.  It is a rule for comparing local frames
at neighboring points.  The symbol $\Omega_\mu$ is the UBT internal or
biquaternionic-frame connection.

Christoffel symbols,
\[
 \Gamma^\rho{}_{\mu\nu},
\]
are closely related, but they act on coordinate indices.  For example,
\[
 \nabla_\mu V^\rho=\partial_\mu V^\rho
 +\Gamma^\rho{}_{\mu\sigma}V^\sigma.
\]
Thus:
\begin{itemize}
 \item $\Gamma^\rho{}_{\mu\nu}$ tells us how the coordinate basis changes;
 \item $\omega_\mu{}^a{}_b$, and its spin/biquaternionic lift $\Omega_\mu$,
       tell us how the local Lorentz frame changes.
\end{itemize}
They are related by the tetrad compatibility identity
\[
 \partial_\mu e_\nu{}^a-\Gamma^\rho{}_{\mu\nu}e_\rho{}^a
 +\omega_\mu{}^a{}_b e_\nu{}^b=0.
\]
They are therefore two descriptions of the same geometric transport, not the
same coefficients.  In particular, $\Gamma=\operatorname{Re}\Omega$ is not a
valid identification.

\section{Why the classical connection is not an extra free field}
For a nondegenerate tetrad, metric compatibility and zero torsion determine a
unique Lorentz connection:
\[
 \boxed{
 \mathring\omega_\mu{}^a{}_b
 =e_b{}^\nu\left(\mathring\Gamma^\rho{}_{\mu\nu}(g)e_\rho{}^a
 -\partial_\mu e_\nu{}^a\right).}
\]
The circle marks the torsion-free Levi--Civita connection.  Thus, in the
ordinary torsion-free GR branch, the frame connection is computed from the
tetrad just as the Christoffel symbols are computed from the metric.  A local
Lorentz rotation changes the numerical coefficients of the tetrad and the
connection but does not change the geometry.

\subsection{Torsion and contorsion}
The torsion two-form is
\[
 T^a=de^a+\omega^a{}_b\wedge e^b.
\]
A general metric-compatible connection can be written
\[
 \boxed{\omega=\mathring\omega(e)+K(T),}
\]
where $K$ is the contorsion.  With
$T^a=\tfrac12T^a{}_{bc}e^b\wedge e^c$ and the convention above,
\[
 \boxed{
 K_{abc}=\frac12\left(T_{cab}-T_{abc}-T_{bca}\right).}
\]
This formula has an important conceptual consequence: once the tetrad and the
torsion are specified, the metric-compatible connection is unique.  The open
full-UBT question is therefore not an arbitrary leftover $\Omega_\mu$; it is
the dynamical law selecting $T$.

If the UBT vacuum equations imply $T=0$, the connection reduces to the
Levi--Civita spin connection.  If spin or another biquaternionic source creates
torsion, the same theorem reconstructs the connection from that torsion.

\section{Flat space and an explicit UBT representer}
In flat Minkowski spacetime, using inertial Cartesian coordinates and a constant
local frame, one can choose
\[
 \Gamma^\rho{}_{\mu\nu}=0,
 \qquad \Omega_\mu=0,
 \qquad D_\mu=\partial_\mu.
\]
This statement can be strengthened from a choice of gauge to an explicit
single-field representer.  Let
\[
 \boxed{
 \Theta_{\mathrm{SR}}(x)=\Theta_0+\sqrt{\mathcal N_0}
 \left(i x^0\mathbf1+x^k\mathbf e_k\right).}
\]
Then
\[
 \mathcal N_0^{-1/2}\partial_0\Theta_{\mathrm{SR}}=i\mathbf1,
 \qquad
 \mathcal N_0^{-1/2}\partial_k\Theta_{\mathrm{SR}}=\mathbf e_k.
\]
The central anticommutator of these four constant directions is the Minkowski
metric.  Every second spacetime derivative of $\Theta_{\mathrm{SR}}$ vanishes,
so it also solves the standard source-free constant-coefficient second-order
master operator in this flat branch.

More generally, for any constant Lorentz tetrad $\bar E_\mu$,
\[
 \Theta_{\mathrm{aff}}(x)=\Theta_0+
 \sqrt{\mathcal N_0}\,\bar E_\mu x^\mu
\]
generates that tetrad exactly.  This closes the special-relativistic
representer problem, although it does not prove uniqueness or curved-space
existence.

In polar coordinates, or in an accelerating frame, $\Gamma$ and $\Omega$ may
be nonzero even though spacetime is flat.  The invariant statement of flatness
is vanishing curvature, not vanishing connection coefficients in every frame.

\section{The four covariant gradients as a tetrad}
Define
\[
 E_\mu:=\frac{1}{\sqrt{\mathcal N_0}}D_\mu\Theta.
\]
The four objects $E_0,E_1,E_2,E_3$ are the local UBT rulers and clock direction.
They are not four extra matter fields; they are the first covariant derivatives
of the single fundamental field $\Theta$.

Let $i$ be the ordinary commuting complex unit and
$\mathbf e_1,\mathbf e_2,\mathbf e_3$ the quaternion units.  In the classical
Lorentz sector write
\[
 E_\mu=i e_\mu{}^0\mathbf1+e_\mu{}^k\mathbf e_k,
 \qquad e_\mu{}^a\in\mathbb R.
\]
Quaternion conjugation changes the signs of the three quaternion units but
leaves the complex scalar coefficient untouched:
\[
 E_\mu^\sharp=i e_\mu{}^0\mathbf1-e_\mu{}^k\mathbf e_k.
\]

\section{Why the anticommutator is the metric}
Quaternion multiplication contains a scalar product and a cross product.  In
the symmetrised product the cross products cancel:
\[
 \boxed{
 \frac12\left(E_\mu^\sharp E_\nu+E_\nu^\sharp E_\mu\right)
 =g_{\mu\nu}\mathbf1.}
\]
The right-hand side contains $\mathbf1$, the unit biquaternion.  The entire
left-hand side is already central: an ordinary real number multiplied by the
algebra unit.  No trace, real-part map, phase projection, or fiber average is
required.

Expanding the coefficients gives
\[
 g_{\mu\nu}
 =-e_\mu{}^0e_\nu{}^0+
  e_\mu{}^1e_\nu{}^1+
  e_\mu{}^2e_\nu{}^2+
  e_\mu{}^3e_\nu{}^3,
\]
which is exactly
\[
 g_{\mu\nu}=e_\mu{}^a e_\nu{}^b\eta_{ab},
 \qquad \eta_{ab}=\operatorname{diag}(-1,1,1,1).
\]
The minus sign of time comes from $i^2=-1$.

\section{What does the other half of the product do?}
The antisymmetric algebraic part is
\[
 \Sigma_{\mu\nu}
 =\frac12\left(E_\mu^\sharp E_\nu-E_\nu^\sharp E_\mu\right).
\]
It describes an oriented plane or bivector and is a natural carrier for local
rotation, boost, and spin information.  It is not itself the curvature.

Three expressions must be kept separate:
\begin{align*}
 \{E_\mu,E_\nu\}_\sharp
 &\longrightarrow \text{metric},\\
 [E_\mu,E_\nu]_\sharp
 &\longrightarrow \text{algebraic bivector/spin structure},\\
 [D_\mu,D_\nu]
 &\longrightarrow \text{connection curvature or gauge field strength}.
\end{align*}

\section{Curvature from noncommuting derivatives}
For a one-sided connection, the curvature has the familiar form
\[
 \mathcal R_{\mu\nu}
 =\partial_\mu\Omega_\nu-\partial_\nu\Omega_\mu
  +[\Omega_\mu,\Omega_\nu].
\]
The final commutator is essential because biquaternion multiplication is
noncommutative.  Geometrically, curvature measures how a frame changes after
transport around a small closed loop.

The equation $E_\mu=D_\mu\Theta/\sqrt{\mathcal N_0}$ creates an additional
integrability condition.  This is the point at which the representation of the
connection becomes decisive.

\section{Why the old rank problem disappears locally}
A symmetric $4\times4$ metric has ten independent components.  Looking only at
one value of $\Theta\in\mathbb B\simeq\mathbb R^8$ suggests the misleading
comparison $8<10$.

The metric is not built from the value of $\Theta$ alone.  It is built from four
covariant derivatives $E_\mu$.  In the Lorentz sector these form a real
$4\times4$ tetrad $e_\mu{}^a$ with sixteen components.  Six changes are local
rotations and boosts that leave the metric unchanged.  Therefore
\[
 16-6=10.
\]
More rigorously, every small symmetric metric variation $h_{\mu\nu}$ is
produced by
\[
 \delta e_\mu{}^a=\frac12 h_{\mu\rho}e^{\rho a}.
\]
The tetrad-to-metric map therefore has rank ten at every nondegenerate tetrad.
This resolves the local kinematic rank question.  It does not yet prove that
every curved tetrad is generated by an on-shell UBT field.

\section{The integrability problem}
Once the classical connection is reconstructed from the tetrad and torsion,
the defining equation becomes schematically
\[
 E_\mu=\mathcal N_0^{-1/2}
 \left(\partial_\mu\Theta+\rho_*(\Omega_\mu[E,T])\Theta\right).
\]
The unknown $E$ occurs on both sides.  This is an implicit nonlinear first-order
PDE or fixed-point system.  Eliminating $\Omega$ kinematically does not by
itself prove existence or uniqueness of a pair $(\Theta,E)$.

\subsection{A one-sided flatness obstruction}
Suppose the connection acts only from the left,
\[
 D^L_\mu\Theta=\partial_\mu\Theta+A_\mu\Theta,
\]
so that
\[
 [D^L_\mu,D^L_\nu]\Theta=F^L_{\mu\nu}\Theta.
\]
In a torsion-free compatible tetrad branch, antisymmetrising the covariant
Hessian gives
\[
 F^L_{\mu\nu}\Theta=0.
\]
If $\Theta$ is invertible as a biquaternion, then
\[
 F^L_{\mu\nu}=0.
\]
Thus the naive one-sided regular representation cannot generically support
curved GR while retaining an invertible $\Theta$ and torsion-free tetrad
compatibility.  Noninvertible stabilizer configurations may evade this result,
but they cannot be assumed to represent arbitrary geometries.

\subsection{The natural two-sided biquaternionic derivative}
Biquaternions naturally admit both left and right multiplication.  Consider
\[
 \boxed{
 D_\mu\Theta=\partial_\mu\Theta+A_\mu\Theta-\Theta B_\mu.}
\]
A direct expansion gives
\[
 \boxed{
 [D_\mu,D_\nu]\Theta
 =F^A_{\mu\nu}\Theta-\Theta F^B_{\mu\nu}.}
\]
Torsion-free compatibility therefore requires
\[
 F^A_{\mu\nu}\Theta=\Theta F^B_{\mu\nu}.
\]
For invertible $\Theta$,
\[
 \boxed{
 F^A_{\mu\nu}=\Theta F^B_{\mu\nu}\Theta^{-1}.}
\]
Unlike the one-sided case, neither curvature must vanish.  The field $\Theta$
intertwines the left and right curvature representations.  This is a genuine
structural narrowing of GAP-10I: the two-sided/bimodule derivative is the
minimal algebra-native route that avoids the generic flatness obstruction.

It is not yet a complete curved-space solution.  The action must still fix the
relation between $A_\mu$ and $B_\mu$, their involution properties, torsion, and
boundary conditions.

\section{Implicit versus transcendental equations}
The words \emph{implicit} and \emph{transcendental} describe different
properties.

An equation is implicit when the unknown occurs inside the map that is supposed
to determine it:
\[
 E=\mathcal F(E,\Theta).
\]
An equation is transcendental when it contains nonalgebraic functions such as
an exponential, logarithm, trigonometric function, or Jacobi theta function:
\[
 E=e^{-E}.
\]
The UBT tetrad equation is already implicit and differential.  If the allowed
field $\Theta(q,\tau)$ is restricted to a Jacobi-theta or another transcendental
function class, then the complete system may be both implicit and
transcendental.  In a formal paper these properties should be stated
separately, even though intuitively both make the unknown ``entangled with
itself.''

\section{Derivatives versus variations}
A derivative such as $D_\mu\Theta$ describes how one field configuration changes
through spacetime.  A variation compares nearby possible configurations:
\[
 \Theta_\varepsilon=\Theta+\varepsilon\,\delta\Theta.
\]
The variation is
\[
 \delta\Theta=
 \left.\frac{d\Theta_\varepsilon}{d\varepsilon}\right|_{0}.
\]
Variations are needed when deriving field equations from an action.  They are
not replacements for spacetime derivatives.  When the connection depends on
the field,
\[
 \delta(D_\mu\Theta)
 =D_\mu(\delta\Theta)+\rho_*(\delta\Omega_\mu)\Theta
\]
and the induced connection variation must be included.

\section{What is closed and what remains open}
The following results are closed:
\begin{enumerate}
 \item the central anticommutator metric;
 \item tetrad-to-metric rank ten with six Lorentz gauge directions;
 \item GAP-10$\Omega$-KIN: unique metric-compatible connection for specified
       tetrad and torsion;
 \item GAP-10$\Omega$-GR: the Levi--Civita spin connection in the torsion-free
       classical branch;
 \item GAP-10L-CONN: preservation of the Lorentz metric and Lorentz slice by a
       metric-compatible connection;
 \item GAP-10I-SR: an explicit affine $\Theta$ representer for Minkowski and
       every constant Lorentz tetrad;
 \item GAP-10I-1S: the naive one-sided invertible curved route is closed as a
       no-go under torsion-free compatibility.
\end{enumerate}
The two-sided derivative narrows but does not completely close the curved
integrability problem.  The following remain open:
\begin{enumerate}
 \item GAP-10T-DYN: derive vacuum torsion or spin-sourced torsion from the UBT
       action;
 \item fix the exact paired left/right representation and involution on
       $\Theta$;
 \item GAP-10I-CURVED: prove local existence, uniqueness, regularity, and global
       continuation for required curved tetrads;
 \item GAP-10L-DYN: prove that the full $\Theta$ dynamics and sources preserve
       the Lorentz slice;
 \item GAP-10D: derive Einstein dynamics rather than only tetrad kinematics;
 \item show that Schwarzschild, Kerr, FRW, and wave sectors are selected by the
       canonical UBT equations.
\end{enumerate}
The earlier compact-$\psi$ fiber closure remains a mathematically interesting
candidate completion, but it is not part of the minimal canonical route.
